\subsection{Finite binary extensions, trace, and Artin--Schreier masks}\label{sec:binary:trace}

For the rest of the chapter, specialize to
$L=\F_{2^n}$ with $n=2^\ell$ and $\ell\geq1$.
The statements that mention uniformity use the finite uniform
distribution from Definition~\ref{def:security:distribution}.
We write $q:=|L|=2^n$.

\begin{lemma}[Frobenius has period $n$]\label{lem:binary:frob-period}
For every $a\in L$, $a^{2^n}=a$.
\end{lemma}
\begin{proof}
If $a=0$, the identity is immediate. Suppose $a\neq0$.
Multiplication by $a$ is a bijection of $L^\times$,
whose cardinality is $q-1$. Define
$P:=\prod_{x\in L^\times}x$, a nonzero element of $L$.
Reindexing the product under the bijection $x\mapsto ax$ gives
\[
 P=\prod_{x\in L^\times}(ax)
  =a^{q-1}\prod_{x\in L^\times}x=a^{q-1}P.
\]
Cancel the nonzero element $P$, obtaining $a^{q-1}=1$.
Multiplying by $a$ yields $a^q=a$, with $q=2^n$.
\end{proof}

\begin{definition}[Absolute trace]\label{def:binary:trace}
For $x\in L$ define
\[
 \Tr(x):=\sum_{j=0}^{n-1}x^{2^j}.
\]
Initially the sum is taken in $L$; the next lemma proves its
value lies in the prime subfield $\F_2$.
\end{definition}

\begin{lemma}[Trace is additive and Frobenius-invariant]\label{lem:binary:trace-add}
For every $x,y\in L$,
\[
 \Tr(x+y)=\Tr(x)+\Tr(y),\qquad
 \Tr(x^2)=\Tr(x).
\]
In particular, $\Tr$ is $\F_2$-linear.
\end{lemma}
\begin{proof}
The first equality follows by applying
Lemma~\ref{lem:binary:frob-add} to each term of the finite sum.
For the second, write
\[
 \Tr(x^2)=x^2+x^{2^2}+\cdots+x^{2^n}.
\]
By Lemma~\ref{lem:binary:frob-period}, the last term is $x$.
Thus the sum is a cyclic permutation of the terms in $\Tr(x)$.
Linearity for the scalar $0$ is immediate, and for scalar $1$
it is the identity $\Tr(1x)=\Tr(x)$.
\end{proof}

\begin{lemma}[Trace takes binary values]\label{lem:binary:trace-binary}
For every $x\in L$, $\Tr(x)\in\F_2=\{0,1\}$.
\end{lemma}
\begin{proof}
By Lemma~\ref{lem:binary:trace-add},
$\Tr(x)^2=\Tr(x^2)=\Tr(x)$.
Hence $\Tr(x)$ is a root of $T^2+T=T(T+1)$.
A field has no zero divisors, so the only possibilities are
$T=0$ and $T=1$, both in $\F_2$.
\end{proof}

\begin{lemma}[Trace is surjective]\label{lem:binary:trace-surj}
The linear map $\Tr:L\to\F_2$ is nonzero and therefore surjective.
\end{lemma}
\begin{proof}
Consider the polynomial
$H(X)=X+X^2+\cdots+X^{2^{n-1}}\in L[X]$.
It is nonzero: its coefficient of $X$ is $1$, and all higher
exponents are distinct. Its degree $2^{n-1}$ is strictly less than
$q=2^n$. If $H(x)=0$ for all $x\in L$, this nonzero polynomial
would have $q$ distinct roots, contradicting
Lemma~\ref{lem:history:rootcount}. Thus some $x$ satisfies
$\Tr(x)=H(x)\neq0$. By Lemma~\ref{lem:binary:trace-binary},
that value equals $1$. Therefore the image of $\Tr$ contains
both $0$ and $1$.
\end{proof}

\begin{corollary}[Sizes of the two trace fibers]\label{cor:binary:trace-fibers}
The sets $\ker\Tr$ and $\{x\in L:\Tr(x)=1\}$ each have
$q/2$ elements.
\end{corollary}
\begin{proof}
Choose $a\in L$ with $\Tr(a)=1$ by
Lemma~\ref{lem:binary:trace-surj}. Translation $x\mapsto x+a$
is a bijection and, by trace additivity, sends elements of trace
zero exactly to elements of trace one and vice versa.
The two fibers partition $L$, and hence have equal cardinality
$q/2$.
\end{proof}

\begin{definition}[Artin--Schreier operator]\label{def:binary:artin-schreier}
Define $\wp:L\to L$ by
\[
 \wp(t):=t^2+t.
\]
\end{definition}

\begin{lemma}[Artin--Schreier is binary-linear]\label{lem:binary:as-linear}
For all $s,t\in L$, $\wp(s+t)=\wp(s)+\wp(t)$.
\end{lemma}
\begin{proof}
Expand $(s+t)^2+s+t$ and apply
Lemma~\ref{lem:binary:square}. The result is
$(s^2+s)+(t^2+t)$.
\end{proof}

\begin{lemma}[Artin--Schreier kernel]\label{lem:binary:as-kernel}
$\ker\wp=\{0,1\}$.
\end{lemma}
\begin{proof}
$\wp(t)=t^2+t=t(t+1)$. This product is zero exactly when
$t=0$ or $t+1=0$, the latter being $t=1$ in characteristic two.
Both values satisfy the equation.
\end{proof}

\begin{lemma}[Every Artin--Schreier value has trace zero]\label{lem:binary:as-trace}
For every $t\in L$, $\Tr(\wp(t))=0$.
\end{lemma}
\begin{proof}
By trace additivity and Frobenius invariance,
$\Tr(t^2+t)=\Tr(t^2)+\Tr(t)=\Tr(t)+\Tr(t)=0$.
\end{proof}

\begin{theorem}[Artin--Schreier characterizes trace zero]\label{thm:binary:as-image}
\[
 \boxed{\operatorname{im}\wp=\ker\Tr.}
\]
Every element of $\ker\Tr$ has exactly two preimages under $\wp$.
\end{theorem}
\begin{proof}
Lemma~\ref{lem:binary:as-trace} gives
$\operatorname{im}\wp\subseteq\ker\Tr$.
By Lemma~\ref{lem:binary:as-linear},
the nonempty fiber above any $y=\wp(t)$ is the coset
$t+\ker\wp$: indeed $\wp(s)=\wp(t)$ is equivalent to
$\wp(s+t)=0$ and hence $s+t\in\ker\wp$.
By Lemma~\ref{lem:binary:as-kernel}, every nonempty fiber
contains exactly two elements. Thus
$|\operatorname{im}\wp|=q/2$.
Corollary~\ref{cor:binary:trace-fibers} shows
$|\ker\Tr|=q/2$. The inclusion between equal-sized finite
sets is an equality.
\end{proof}

\begin{corollary}[Uniform trace-zero masking]\label{cor:binary:as-uniform}
If $T$ is uniform on $L$, then $\wp(T)$ is uniform on
$\ker\Tr$.
\end{corollary}
\begin{proof}
By Theorem~\ref{thm:binary:as-image}, every trace-zero
element has exactly two preimages. Therefore its probability
as the output $\wp(T)$ is $2/q=1/(q/2)$, independent of
the chosen trace-zero element.
\end{proof}

\begin{theorem}[Perfect hiding of an individual trace-zero shift]\label{thm:binary:single-mask}
Fix $v\in\ker\Tr$ and let $T$ be uniform on $L$.
Then
\[
 U:=v+\wp(T)
\]
is uniform on $\ker\Tr$, independently of $v$.
\end{theorem}
\begin{proof}
Corollary~\ref{cor:binary:as-uniform} makes
$\wp(T)$ uniform on the additive subgroup $\ker\Tr$.
Since $v$ lies in that subgroup, translation
$y\mapsto v+y$ is a bijection of $\ker\Tr$,
so $U$ is uniform there. Its distribution is therefore
independent of the fixed value $v$.
\end{proof}

\begin{remark}[No joint-ZK conclusion]
Theorem~\ref{thm:binary:single-mask} concerns one public
field element. If the same $T$ is subsequently used in
another published equation, the joint transcript can have
correlations. The complete verification protocol must
analyze those correlations, as illustrated by the
reused-mask attack of Theorem~\ref{thm:kbff:one-query-attack}.
\end{remark}
